\clearpage
\part*{Part XV. Quantum Information}
\label{atlaspart:15-quantum-information}
\addcontentsline{toc}{part}{Part XV. Quantum Information}
\section*{Introduction to Part XV}\label{narr:part:XIII}
\addcontentsline{toc}{section}{Introduction to Part XV}
The construction began as quantum information, and Part XV returns it there. The closure of a web under the moves is a computation, the moves are a gate set, the coherent read is the gate the literature calls the Hadamard, the read is the maximal distillation, the enumeration of the records is the cost of a read, and the transport of a basis around a closed path has a holonomy. The Part is the one where the vocabulary of the axiom and the vocabulary of the literature come closest, since both speak of qudits, entangled pairs and gates, and the reader will find more theorems here than in any Part beyond the framework.

The order of the six chapters is the order of a course on quantum computation read against the axiom. The first chapter identifies the objects, the bit as the pair ket, the gate as the act at one link and the channel as the wash, and shows that every ready name of the computation vocabulary is a corollary or a gloss. The second and third chapters ask what the gate set is and whether the one gate the literature cannot do without, the Hadamard, lies inside it. The fourth chapter reads distillation, which is the act of the axiom seen as an operation on pairs, and the fifth reads the cost of a read as a count of records and the adiabatic theorem as a statement about the records a move does not touch. The sixth chapter reads the transport of a basis along a path of records and the holonomy around a closed path, which is the object Parts IV, VII and VI return to. A reader who wants the one result that ties the Part to the axiom should read the first theorem of chapter 13-4, where the distilling line of the text is shown to be a corollary of the one word maximally entangled.

Chapter 13-1 reads the closure as the computation in Theorem~230.1, a theorem resting on the axiom chapter, the records, the moves, the saturation and the probe. Every ready name of the computation vocabulary is a corollary of the text or of a theorem, or a gloss named beside an object the construction has, and every gated one names its gate, the bit being the qubit pair ket of the axiom, the gate the act at one link and the channel the wash at one act, Theorem~230.1. Chapter 13-2 states the gate set at the Weyl criterion, Open result~231.1, as an open result with its parting named: the five finite groups of the gate-set result, the one-register gates, the step transport a permutation and the holonomy a phase at the roots of unity of the ring, are re-met from its tables as exact matrix groups at the orders two, six, one, two and 2304 and compared as abstract groups with the criterion, where the parting is named, Open result~231.1. Chapter 13-3 reads the Hadamard as the coherent read, Open result~232.1, open, with the Hadamard inside the finite gate group at the pair register, Theorem~232.2, a theorem. The coherent read of the qubit-qubit link is written as an operator at thirty-two candidates from the record spectrum of the count law, Open result~232.1, the candidate read from the readings table is the Hadamard reached at the invariant cut at order two, and the Hadamard tensor the identity is a member of the finite gate group at the pair register with the order 2304 unchanged, so that no move beyond the twenty-eight is admitted, Theorem~232.2. The reader who expects a universal gate set should read the two boxes together: the construction has a finite gate group, the Hadamard lies inside it, and the question of the Weyl criterion is the question of what that finite group is as an abstract group.

Chapter 13-4 is the distillation chapter. Theorem~233.1 obtains the distilling line as a corollary of the axiom, Theorem~233.2 reads the read as the maximal distillation, Theorem~233.3 reads every census move as one act of distilling, and Theorem~233.4 reads the subdivision as a path of four one-act moves, three theorems and one conditional result. The distilling line of the text is a corollary of the axiom and of the qudit line, since the output of an act exists only as a configuration of maximally entangled pairs and the act on two qudits is a projection onto a pair on the common support, the rank identity holding at sixty-nine of sixty-nine operators, Theorem~233.1. The first clause of the read line is a washed clause, a corollary at the level of the known, the basis-matched restriction of the pair projections being the family of maximally entangled projections characterized by the certainty the known clause names, Theorem~233.2. Conditional on the one-act clause, re-met at 32851 of 32851 census instances, the distilling projection makes the two joined directions a pair and composes their two external pairs into one pair of the far ends, Theorem~233.3, and the re-typing of the 1653 subdivisions of the census as paths of four one-act moves is forced by the text at the pair-ket unit, every subdivision decomposing into one read and three joins, Theorem~233.4. Chapter 13-5 reads the cost: Theorem~234.1 reads the enumeration of the records as the cost of a read and Theorem~234.3 is the adiabatic theorem at the untouched records, both theorems, with the class partition under one move, Open result~234.2, open between them. At the sixteen networks of the item and at the sixty-four library networks the records of every link are counted by the closed count law at the slot polarity by order zero to three and by kind, the orders and the kinds summing to the records at every link and network, Theorem~234.1. The twenty-eight moves applied once at every library network give 5707 applications, the class partition by genre and kind carried whole at 4124 of them and parting at 1583, the first parting the attach at the odd leg of the first library network, Open result~234.2, while at every one of those applications every class of the records at the links touching no move site lands whole in the class of the same name at the result, so that the adiabatic theorem is derived at the untouched records with no parting, Theorem~234.3.

Chapter 13-6 reads the transport. Theorem~235.1 reads the read transport from the pair, Theorem~235.2 the holonomy of the basis transport around a closed path, Theorem~235.3 the step generator as a number operator, and Theorem~235.4 the carrier phase datum at the four-cycle, four theorems, with the four open results beside the transport gathered in Open result~235.5. The far set is an orthogonal basis of the far support at every genre and orientation, and the orientation sign is needed at no read, Theorem~235.1. At the 341 closed paths of the deciding set of the library the transport composed around the path at the identity matching is an exact unitary on the register of the starting end, defined in the ring of the twenty-fourth root of unity, nontrivial at the 333 odd paths and trivial at the eight even ones, Theorem~235.2. At every one of the 1692 registers of the step-generator result the step generator equals the ring quotient multiplied by the cycle-index operator, a number operator with an equally spaced spectrum whose spacing is the quotient of that register and whose level count is the cycle order, Theorem~235.3, and at every carrier carrying a qutrit-qubit register the carrier phase datum equals the primitive fourth roots of unity in the holonomy spectrum of the transport-generator four-cycle, present at the W, the Z and the photon words and absent at the gluon, Theorem~235.4. The hyperfine interval of chapter 5-4 and the discrete laws of chapter 3-6 rest on results of this Part, which is why the derivation graph shows edges running from Part XV back into Parts IV and VII.

The adiabatic theorem deserves one more sentence, since its name carries a continuum in the literature. Here it is a statement about a partition: a move touches a set of sites, the records at the links that touch none of those sites are partitioned by class, and the theorem says that the partition is carried whole across the move at every one of the 5707 applications, Theorem~234.3. The condition of the literature, a gap that is not crossed while a parameter varies slowly, has no counterpart, because nothing varies slowly; the move is one act, and the records it does not touch are counted on both sides of it. The same finiteness governs the cost of a read, which is an enumeration of records by order and kind rather than a count of gates in a circuit, Theorem~234.1.

In the language of the literature, this Part names computation, the gate set, the Hadamard gate, the Weyl criterion, distillation, the adiabatic theorem, holonomy and the number operator. The correspondence is by name at each. The boxes establish a finite gate group, a maximal distillation and a holonomy around a closed path of records.

The dependency lines of the Part reach back to the axiom chapter, the records, the moves, the saturation and the probe of Part I, and forward from this Part into Parts IV, VII, VI and XIV, which read the holonomy, the carrier phase datum and the step generator by number. The gate set rests on the moves, the probe and the closure, the Hadamard on the gate set and the count law, distillation on the axiom and the one-act clause, and the transport on the pair and the read line. A reader following the derivation graph will find that Part XV is cited from more Parts than it cites, which is the signature of a Part that reads the axiom itself rather than the objects built from it, and will find at its first chapter the shortest statement in the compendium of what the construction is: a finite set of maximally entangled pairs, the acts that distill them, and the records the acts leave.
\section{The closure as the computation}\label{mod:XIII-1-closure-as-the-computation}
\chdesig{15-1}{13-1}{XIII-1-closure-as-the-computation}{XIII-1-closure-as-the-computation}
\label{narr:open:13-1}The closure is a computation. Theorem~230.1 reads it so, a theorem resting on the axiom chapter, the records, the moves, the saturation and the probe.
\begin{chtheorem}[The closure as the computation]\label{box:XIII-1-closure-as-the-computation:1}
Every ready name of the computation vocabulary is a corollary of the text or of a theorem, or a gloss named beside an object the construction has, and every gated one names its gate. The bit is the qubit pair ket of the Axiom; the gate is the act at one link; the channel is the wash at one act; the measurement is the wash at a cap; the classical register is the washed record; teleportation is the read passed along the pair together with its record; the Schmidt rank is the common support and the purification the Born face; the cost of a computation is the record count of a read; the circuit is the closure of the act from a root, with the saturated networks its halting states and the stop condition of the nuclear roster its halting rule; and the Turing machine stands open at the universality question. Of the thirty-six names (fifteen ready, twenty-one gated) nine are derived, fourteen glossed and thirteen open at their named gates (the Turing machine and its thesis, the halting problem, the Kolmogorov complexity, the algorithms, the codes, the key distribution, the demon, the accessible-record bound and the adiabatic form); the cost of a read is printed at 18 of 18 networks with the sum equal to the orders and at twelve closures (six stopped, six at the web cap). The words bit, gate, circuit, channel, teleportation, Turing machine and complexity class enter by name.
\end{chtheorem}
\chstanding{theorem}
\chdeps{Definitions 1.1 to 1.6, Theorems 2.1 to 2.3, Open results 2.4 and 2.5, Theorem 2.6, Open results 2.7 and 2.8 and Theorems 2.9 and 2.10, Theorems 18.1 to 18.4 and Open results 18.5 and 18.6, Theorems 28.1 and 28.2, Open results 28.3 and 28.4 and Theorem 28.5, Open results 241.1 and 241.2, Theorems 241.3 and 241.4 and Open results 241.5 to 241.7 and Open results 245.1 and 245.2, Theorem 245.3, Open results 245.4 to 245.8 and Theorem 245.9}{Open result 231.1 and Theorem 234.1, Open result 234.2 and Theorem 234.3}
The theorem is conditional on the text, on the open results it rests on, on the Lean package of this compendium and on the complete list of computation names and the roster library. What stands open is the thirteen names typed open at their gates, the Turing machine among them, which the construction meets at the gate set of the moves and at the Hadamard in the two chapters that follow.
\begin{chbridge}
The literature calls the pair ket a Bell state (exact), the act a Bell-basis projection, an entanglement swap (exact), the wash the twirl (exact), the closure the forward closure of a set under operations (exact), the read a projective measurement passed along the pair (nearest), and the objects named the circuit model of a quantum computation \cite{Nielsen2000}, the Gottesman-Knill theorem \cite{Gottesman1998} and the measurement-based model \cite{Raussendorf2001}, by name only.
\end{chbridge}
\begin{chproof}[finite computation and census]
The premises are the Axiom and the six definitions of Section~\ref{mod:FRAMEWORK} (definitions), the axiom chapter of Section~\ref{mod:I-1-axiom-and-the-six-definitions} (an open result, read at its face), the records and the Born count of Section~\ref{mod:I-3-records-reads-and-the-born-count-the-gain-lemma}, the moves of Section~\ref{mod:I-4-moves-in-one-table}, the saturation of Section~\ref{mod:I-5-saturation-in-one-figure-seeds-waves-termination-sinks} and the read at a register of Section~\ref{mod:I-6-read-at-a-register-the-probe-and-the-measurement} (theorems). Nothing outside that list is used, and the chain from the axiom to this result is five chapters deep.
\emph{The census part.} The set counted is the thirty-six computation names of the literature not already graded derived, fifteen ready and twenty-one gated. From a name, by its map to a root of the construction (a line of the text, a theorem or an object) and the standing of that root, we obtain its type by the typing rule. Derived at a proved root with an exact check, glossed at an established object with the missing clause named, open at a gate named. The counts by class are
\begin{equation}
\text{derived } 9, \qquad \text{glossed } 14, \qquad \text{open } 13, \qquad 9 + 14 + 13 = 36 ,
\end{equation}
with the twelve gate names, the fourteen circuit names and the ten cost names each at their claims, every ready name at its rule (15 of 15), every gated name at its gate (21 of 21) and the first name typed at none none. The smallest cases. The bit, a qudit of dimension two at a maximally entangled pair of the Axiom read in one of its bases fixing one of two directions at its far end, derived at three roots of three. The circuit, the closure of the act from a root with the saturated networks its halting states, glossed at four roots with the name laid beside the closure object. The Turing machine, open at the gate of the cost of a computation and the universality question.
\emph{The finite computation.} The inputs are the sixteen networks of the item and the twelve closures of the fixed library with their orders tables. The procedure is this. For every network, count the records of every link by the closed count law at the slot polarity, order by order from zero to three, and sum the orders. The cost of a read at a network at an order is the record count at that order. At the electron body,
\begin{equation}
1701 + 1854 + 130 + 0 = 3685 ,
\end{equation}
the orders zero to three summing to the records, and at the muon body $2637 + 2868 + 170 + 0 = 5675$, at the tau $3573 + 3882 + 210 + 0 = 7665$, at the docked up quark $158 + 180 + 6 + 0 = 344$, the sum equal to the orders at 18 of 18 networks. At the twelve closures the rewriting letters are attach, hub, join and pair with the root at the order zero, six closures stopped (the smallest at one order with four webs, three saturated) and six unstopped at the web cap at the fifth order with 504 webs, the growth order the step count by name and the web count the space by name. The result is the theorem. Every ready name derived or glossed at a verified root, every gated name at its gate, the cost of a read a count at every network.
The check is by the counter-tests, a root struck moves 73 of 496 cells and every deciding cell. The closure read at zero acts moves 29 cells and the circuit cell. The text struck moves 491 cells and withholds every cell. Ten controls hold, the partition of the computation names exact and the per-link count cells summing to the orders at every network among them.
The result was reproduced by an independent route, two engines written apart agreeing at every name, root and count.
\chcert{36}
\end{chproof}

\par\medskip\noindent{\itshape The computation named, the next chapter reads its gate set.}\label{narr:close:13-1}\par
\section{The moves as a gate set and the Weyl criterion}\label{mod:XIII-2-moves-as-a-gate-set-and-the-weyl-criterion}
\chdesig{15-2}{13-2}{XIII-2-moves-as-a-gate-set-and-the-weyl-criterion}{XIII-2-moves-as-a-gate-set-and-the-weyl-criterion}
\label{narr:open:13-2}Open result~231.1 states the gate set at the Weyl criterion as an open result with its parting named, resting on the moves, the probe and the closure.
\begin{chwall}[The gate set at the Weyl criterion]\label{box:XIII-2-moves-as-a-gate-set-and-the-weyl-criterion:1}
The five finite groups of the gate-set result (the one-register gates, the step transport a permutation and the holonomy a phase at the roots of unity of the ring) are re-met from its tables as exact matrix groups at the orders two, six, one, two and 2304, and compared as abstract groups with the Weyl groups of the stabiliser algebras of the construction (the symmetric group of two letters at the qubit register, of three at the qutrit, of four at the hub register of dimension four, the product of two four-letter groups at the two ququart ends of the dock register, the product of two two-letter groups at the pairing register): equal at one of five registers, the qubit, where the move group of order two maps onto the two-letter group by its generator to the reflection, and parting at four of five. The first parting register is the qutrit, where the move group of order six (a label transposition multiplied by a holonomy phase of order three, abelian and cyclic) parts from the symmetric group of three letters at the relations with the orders equal; beyond it the ququart register parts at the order (trivial against twenty-four), the dock register at the order (two against 576) and the pairing register with the move group larger (2304 against four, the move group outside the normaliser of every torus). The Hadamard is an element of order two outside the normaliser of the torus and outside the move group at the qubit register, reached by no move, while at the pairing register the Hadamard tensor the identity stands inside the move group; universality is taken up in the next chapter, and the circuit name is spelled as a move sequence with its gap named.
Re-posed at the first parting register: the four holonomy letters at the qutrit start, re-met as exact monomials of the ring of the twenty-fourth root, are each the one transposition of the labels one and two at a phase or the identity and generate the cyclic group of order six with one involution, while the Weyl group of su(3) at that register is the symmetric group of three letters with three reflections and six relations holding at six of six; the relations checked one by one at the gate set in their order leave six of the thirty-six ordered pairs at the first, a reflection, and none at the second, a second reflection distinct from the first, so the gate set at the qutrit register does not meet the Weyl criterion at every relation and parts from the Weyl group of su(3) at the second relation, the readings beside parting at the sixth relation and at the second.
\end{chwall}
\chstanding{open result}
\chdeps{Definitions 1.1 to 1.6, Theorems 28.1 and 28.2, Open results 28.3 and 28.4 and Theorem 28.5, Theorem 230.1 and Open results 241.1 and 241.2, Theorems 241.3 and 241.4 and Open results 241.5 to 241.7}{Open result 60.1, Theorem 93.1 and Open result 232.1 and Theorem 232.2}
At one face the move group would equal the Weyl group of the stabilizer algebra at every register. At the face this chapter inherits it is equal at the qubit register alone and parts at the four beyond, the first parting at the qutrit named with both objects below. The universality of the moves as a gate set is a question met at the Hadamard in the next chapter.
\begin{chbridge}
The literature calls the register the tensor factors of the common support (nearest), the holonomy the parallel transport around a closed path (nearest), the move an elementary operation on the network of pairs (nearest), the step transport a permutation of the register basis (nearest), and the objects named a universal gate set, the Clifford group as the normalizer of the Pauli group \cite{Gottesman1998,Nielsen2000} and the Weyl group of a Lie algebra as the reflections of its root system \cite{Weyl1927}, by name only.
\end{chbridge}
\begin{chproof}[census and exact cell]
The premises are the Axiom and the six definitions of Section~\ref{mod:FRAMEWORK} (definitions), the moves of Section~\ref{mod:I-4-moves-in-one-table} and the read at a register of Section~\ref{mod:I-6-read-at-a-register-the-probe-and-the-measurement} (theorems) and the closure as the computation of Section~\ref{mod:XIII-1-closure-as-the-computation} (a theorem). Nothing outside that list is used, and the chain from the axiom to this result is six chapters deep.
\emph{The exact part.} The object is the move group at the qubit register, generated by the holonomies of the hub triangle at its qubit start, the identity and the letter $XZ$. From the two generators, by closure under multiplication over the ring of the twenty-fourth root of unity $z$, we obtain the group
\begin{equation}
G_{2} = \{ \mathbf{1}, \ \mathrm{diag}(1, z^{12}) \} = \{ \mathbf{1}, \ \mathrm{diag}(1, -1) \}, \qquad |G_{2}| = 2 ,
\end{equation}
cyclic, every element diagonal, inside the torus, with one permutation matrix of two. From the stabilizer algebra of the register at the reading of reference (the span of dimension three at rank one, named $\mathfrak{su}(2)$ by the literature), by the adjacent transpositions at the Coxeter relations, we obtain its Weyl group, the symmetric group of two letters of order two. From the two groups, by the explicit map of the first generator to the identity and the second to the reflection, we obtain the isomorphism, so the qubit register is equal. At the qutrit register the four letters of the core triangle generate
\begin{equation} \adjustbox{max width=\dimexpr\linewidth-3em\relax}{$\displaystyle \begin{gathered} |G_{3}| = 6, \qquad \text{abelian, the element orders } 1\!:\!1,\ 2\!:\!1,\ 3\!:\!2,\ 6\!:\!2, \\ \text{against } S_{3}: \ 1\!:\!1,\ 2\!:\!3,\ 3\!:\!2, \ \text{non-abelian}, \end{gathered}$} \end{equation}
the orders equal and the relations parting (every pair of generators commuting in $G_{3}$, the Coxeter relation $(s_{1} s_{2})^{3} = 1$ with $s_{1} s_{2} \neq s_{2} s_{1}$ in $S_{3}$), which is the first parting element.
\emph{The census part.} The set counted is the five registers of the gate-set result. From each, by the move group re-met from the tables (the orders two, six, one, two and 2304) and the Weyl group of its stabilizer algebra (two, six, twenty-four, 576 and four), we obtain the comparison as abstract groups, equal at 1 of 5 (the qubit) and parting at 4 of 5. The qutrit at the relations, the ququart at the order (one against twenty-four), the dock at the order (two against 576) and the pairing register with the move group the larger (2304 against four, the move group outside the normalizer of every torus). The move group lies inside the normalizer of the torus at 4 of 5 registers. From the Hadamard at the qubit register, by the test against the group, we obtain an element of order two outside the normalizer of the torus and outside the move group, reached by no move, the first gate with no move, while at the pairing register the Hadamard tensor the identity lies inside the move group of order 2304. The readings beside (the stabilizer read as an orthogonal algebra, the Weyl groups of type $B$ or $D$) part at the order at every register beyond the qubit and are printed, not chosen.
The check is by the counter-tests. The holonomy phases struck at the qutrit register move the move group there to the transposition of order two and the parting to the order. The comparison read at the orders alone moves the first parting register to the ququart. The text struck leaves the gate unmet. The gate set replaced by the Weyl representatives moves the deciding cell to equal at every register. Seven controls were posed and six hold, the one that differs a reality test mis-posed at one engine at one shared key, no identity of the tables, so the control cell reads as printed and the result stands at its face.
The result was reproduced by an independent route, two engines written apart, one exact at the ring and one at complex floats keyed to six places, agreeing at every group and comparison.
\chcert{5}
\end{chproof}
\begin{figure}[h]
\begin{center}
\setlength{\unitlength}{1pt}
\begin{picture}(200,90)(0,0)
\put(40,20){\circle*{6}}
\put(160,20){\circle*{6}}
\put(100,80){\circle{18}}
\put(40,20){\line(1,0){120}}
\put(46,23){\line(1,1){48}}
\put(154,23){\line(-1,1){48}}
\put(30,6){$f_{1}$ (2)}
\put(150,6){$f_{2}$ (2)}
\put(112,76){$F$ (4)}
\put(88,8){$(2,2)$}
\put(40,52){$(2,4)$}
\put(138,52){$(4,2)$}
\end{picture}
\end{center}
\caption{The hub triangle at its qubit start. Two qubit sites $f_{1}$, $f_{2}$ and the hub $F$ of dimension four, the closed path $f_{1} \to f_{2} \to F \to f_{1}$ whose holonomy at the letter $XZ$ generates the move group of order two at the qubit register.}
\end{figure}

\par\medskip\noindent{\itshape One gate of the set is read on its own in the next chapter.}\label{narr:close:13-2}\par
\section{The Hadamard as the coherent read}\label{mod:XIII-3-hadamard-as-the-coherent-read}
\chdesig{15-3}{13-3}{XIII-3-hadamard-as-the-coherent-read}{XIII-3-hadamard-as-the-coherent-read}
\label{narr:open:13-3}Open result~232.1 reads the Hadamard as the coherent read, an open result, and Theorem~232.2, a theorem, places the Hadamard inside the finite gate group at the pair register.
\begin{chwall}[The Hadamard as the coherent read]\label{box:XIII-3-hadamard-as-the-coherent-read:1}
The four moves as gates are re-met at both computations (the five one-register groups at the orders two, six, one, two and 2304), and the coherent read of the qubit-qubit link is written as an operator at thirty-two candidates from the record spectrum of the count law: at the reading of reference the verdict is no gate; at the readings beside the Hadamard is reached at the cut records of the holonomy-invariant subspace up to a twenty-fourth-root phase and a real scale. The group generated at the qubit register by the one-register gates with the coherent read is finite at the move group alone (order two), finite with the Hadamard adjoined (order sixteen, 192 with the phases) and infinite with the whole end-transport reading adjoined; so the computation domain is finite by right at the reading of reference, an open result.
\end{chwall}
\chstanding{open result}
\chdeps{Definitions 1.1 to 1.6, Theorems 18.1 to 18.4 and Open results 18.5 and 18.6, Theorems 28.1 and 28.2, Open results 28.3 and 28.4 and Theorem 28.5 and Open result 231.1}{no theorem beyond it at this release}
At one face the computation domain would be universal by the construction, the coherent read a gate whose group with the one-register gates is infinite. At the face this chapter inherits the coherent read is no gate at the reading of reference and the group at one register is the move group of order two, finite.
\begin{chtheorem}[The Hadamard inside the finite gate group at the pair register]\label{box:XIII-3-hadamard-as-the-coherent-read:2}
The candidate read from the readings table of that result is the Hadamard reached at the invariant cut at order two, and the Hadamard tensor the identity is a member of the finite gate group at the pair register (the order 2304 unchanged), so no move beyond the twenty-eight is admitted, the gate set is closed under the candidate and is the limit of the move set, and the open result is a residue of the finite group; the infinite order belongs to the invariant-cut rotations alone, beyond the register face at every seeded sink carrying a qubit-qubit link (46412 of 46563 sinks), candidates the saturation cannot admit under the termination theorem.
\end{chtheorem}
\chstanding{theorem}
\chdeps{Definitions 1.1 to 1.6, Theorems 18.1 to 18.4 and Open results 18.5 and 18.6, Theorems 28.1 and 28.2, Open results 28.3 and 28.4 and Theorem 28.5 and Open result 231.1}{no theorem beyond it at this release}
The theorem is a composition riding the lowest grade in its chain, conditional on the open results it rests on, on the termination theorem of the saturation and on the three tables of the seeded sinks. It types the open result a residue of the finite group and re-opens the move table at nothing.
\begin{chbridge}
The literature calls the wash the twirl (exact), the register the tensor factors of the common support (nearest), the holonomy the parallel transport around a closed path (nearest), the record space the Hilbert space spanned by the records (exact), and the objects named the Hadamard gate of order two with the phase gate generating the single-qubit Clifford group \cite{Nielsen2000}, the classical simulability of a Clifford circuit \cite{Gottesman1998} and the magic states distilled from noisy copies \cite{BravyiKitaev2005}, by name only. The construction supplies the states, projections and channels it uses at Sections~\ref{mod:FRAMEWORK}, \ref{mod:I-3-records-reads-and-the-born-count-the-gain-lemma} and \ref{mod:II-4-measurement-the-wash-at-one-act-and-the-povm}.
\end{chbridge}
\begin{chproof}[census]
The premises are the Axiom and the six definitions of Section~\ref{mod:FRAMEWORK} (definitions), the records and the Born count of Section~\ref{mod:I-3-records-reads-and-the-born-count-the-gain-lemma} and the read at a register of Section~\ref{mod:I-6-read-at-a-register-the-probe-and-the-measurement} (theorems), and the gate set of Section~\ref{mod:XIII-2-moves-as-a-gate-set-and-the-weyl-criterion} (an open result, read at its face). Nothing outside that list is used, and the chain from the axiom to this result is seven chapters deep.
The set counted is the sixteen readings of the coherent read of the qubit-qubit link with their invariant parts, thirty-two candidates. The four moves as gates are re-met first. The step transport at every genre a permutation matrix of the register (at $(2,2)$ the cycle type $1^{2} 2^{1}$ of order two, the matrix $S_{(2,2)}$ with the rows $(1, 0, 0, 0)$, $(0, 0, 1, 0)$, $(0, 1, 0, 0)$ and $(0, 0, 0, 1)$, equal to the step table at 7 of 7 rows), the holonomy a phase at the roots of unity of the ring at 24 of 24 lines of its table, the pairing an isometry whose image of the product state is the pair ket $(1, 0, 0, 1)$ and the wash a channel whose Kraus operators are the two cap projectors of the near end tensor the identity of the far end. From the $(2,2)$ count law (the center two and the far two, the roster 54, the fixed 6, the records 57 at the orders $0\!:\!24$, $1\!:\!30$, $2\!:\!3$, the registers whole 16 and cut 9), by the explicit slot-letter records, we obtain the operator of the coherent read at each reading. At the reading of reference (the one end register of dimension two at the letters $Z$, $X$, $XZ$, every center letter transported to the far end by the transpose and multiplied with the far letters in slot order) the operator is
\begin{equation} \adjustbox{max width=\dimexpr\linewidth-3em\relax}{$\displaystyle \begin{gathered} M = \begin{pmatrix} 6 + 18 z - 18 z^{5} + 18 z^{3} & -4 z - 4 z^{3} + 4 z^{5} \\ -12 - 4 z - 4 z^{3} + 4 z^{5} & -6 + 10 z + 10 z^{3} - 10 z^{5} \end{pmatrix}, \\ M M^{\dagger} \ \text{not a multiple of the identity}, \end{gathered}$} \end{equation}
so the verdict is no gate, which is the property tested at the reading of reference. At the six readings at the invariant cut (the cut records of the holonomy-invariant subspace of dimension nine) the operator is the Hadamard up to a twenty-fourth-root phase and a real scale. From the one-register gates with a reading adjoined, by closure under multiplication at the caps, we obtain the generated group at the qubit register, the move group alone of order two, finite, the first gate with no move the Hadamard, with the invariant-cut reading adjoined of order sixteen, finite, the Hadamard inside, with the phases adjoined of order 192, finite, and with the whole end-transport reading adjoined infinite. The criterion, stated once, is that the domain is universal by the construction if and only if the coherent read is a gate at the reading of reference and the generated group is infinite. The first conjunct fails, so the domain is finite by right there, the counts by class finite at 3 readings and infinite at 1 among the four groups closed.
The check is by the counter-tests, the far letters struck move 230 cells, the coherent read first. The transpose struck at the transport moves 443 cells. The text struck refuses the gate. The invariant subspace replaced moves 147 cells, the generated group first. Fourteen controls were posed and eleven hold, the three that differ at presentation cells of a parent table, so the control cell reads as printed and the result stands at its face.
The result was reproduced by an independent route, two engines written apart agreeing at every reading and every group.
\chcert{32}
\end{chproof}

\par\medskip\noindent{\itshape The next chapter reads the read itself as a distillation.}\label{narr:close:13-3}\par
\section{Distillation: the read as the maximal distillation and the distillation theorem}\label{mod:XIII-4-distillation-the-read-as-the-maximal-distillation-and-the-distillation-theorem}
\chdesig{15-4}{13-4}{XIII-4-distillation-the-read-as-the-maximal-distillation-and-the-distillation-theorem}{XIII-4-distillation-the-read-as-the-maximal-distillation-and-the-distillation-theorem}
\label{narr:open:13-4}Theorem~233.1 obtains the distilling line as a corollary of the axiom, Theorem~233.2 reads the read as the maximal distillation, Theorem~233.3 reads every census move as one act of distilling, and Theorem~233.4 reads the subdivision as a path of four one-act moves.
\begin{chtheorem}[The distilling line as a corollary of the Axiom]\label{box:XIII-4-distillation-the-read-as-the-maximal-distillation-and-the-distillation-theorem:1}
The distilling line of the text is a corollary of the Axiom and of the qudit line: the output of an act exists only as a configuration of maximally entangled pairs, the act on two qudits is a projection onto a pair on the common support (the rank identity at 69 of 69 operators), and the pair lives on the common support alone; the four steps are derived with no posed step.
\end{chtheorem}
\chstanding{theorem}
\chdeps{Definitions 1.1 to 1.6, Theorems 2.1 to 2.3, Open results 2.4 and 2.5, Theorem 2.6, Open results 2.7 and 2.8 and Theorems 2.9 and 2.10, Theorems 9.1 to 9.7 and Open result 9.8, Theorems 18.1 to 18.4 and Open results 18.5 and 18.6 and Open results 241.1 and 241.2, Theorems 241.3 and 241.4 and Open results 241.5 to 241.7}{no theorem beyond it at this release}
\begin{chtheorem}[The read as the maximal distillation]\label{box:XIII-4-distillation-the-read-as-the-maximal-distillation-and-the-distillation-theorem:2}
The first clause of the read line is a washed clause, a corollary at the level of the known: the basis-matched restriction of the pair projections is the family of maximally entangled projections, characterised by the certainty the known clause names (a read at every basis if and only if the reduced state is flat, the Gram identity at 120 of 120); the basis of that clause is a label beyond the projection, washed at the matching group, and the read cap, the dimension plus one, is the count of mutually unbiased bases.
\end{chtheorem}
\chstanding{theorem}
\chdeps{Definitions 1.1 to 1.6, Theorems 2.1 to 2.3, Open results 2.4 and 2.5, Theorem 2.6, Open results 2.7 and 2.8 and Theorems 2.9 and 2.10, Theorems 9.1 to 9.7 and Open result 9.8, Theorems 18.1 to 18.4 and Open results 18.5 and 18.6 and Open results 241.1 and 241.2, Theorems 241.3 and 241.4 and Open results 241.5 to 241.7}{no theorem beyond it at this release}
\begin{chconditional}[Every census move one act of distilling]\label{box:XIII-4-distillation-the-read-as-the-maximal-distillation-and-the-distillation-theorem:3}
Conditional on the one-act clause (one act per census move, re-met at 32851 of 32851 instances and exceeded by the order rule of the library at 3088), the distilling projection makes the two joined directions a pair and composes their two external pairs into one pair of the far ends, and every census move makes or unmakes the pair kets its typed single act names and no other, the induced matching equal to the persistence face and differing from the order face at the 41 re-matched joins.
\end{chconditional}
\chstanding{conditional theorem}
\chdeps{Definitions 1.1 to 1.6, Theorems 2.1 to 2.3, Open results 2.4 and 2.5, Theorem 2.6, Open results 2.7 and 2.8 and Theorems 2.9 and 2.10, Theorems 9.1 to 9.7 and Open result 9.8, Theorems 18.1 to 18.4 and Open results 18.5 and 18.6 and Open results 241.1 and 241.2, Theorems 241.3 and 241.4 and Open results 241.5 to 241.7}{no theorem beyond it at this release}
\begin{chtheorem}[The subdivision as a path of four one-act moves]\label{box:XIII-4-distillation-the-read-as-the-maximal-distillation-and-the-distillation-theorem:4}
The re-typing of the 1653 subdivisions of the census as paths of four one-act moves is forced by the text at the pair-ket unit of the first reading, every subdivision decomposing into one read and three joins; under that unit the one-act clause is a theorem at every instance, and at the unit of the projection with its common-support byproduct the subdivision is one inverse act at 1448 instances and the path alone at 205, both readings printed and neither chosen.
\end{chtheorem}
\chstanding{theorem}
\chdeps{Definitions 1.1 to 1.6, Theorems 2.1 to 2.3, Open results 2.4 and 2.5, Theorem 2.6, Open results 2.7 and 2.8 and Theorems 2.9 and 2.10, Theorems 9.1 to 9.7 and Open result 9.8, Theorems 18.1 to 18.4 and Open results 18.5 and 18.6 and Open results 241.1 and 241.2, Theorems 241.3 and 241.4 and Open results 241.5 to 241.7}{no theorem beyond it at this release}
Every theorem is conditional on the text and on the open results, the conditional theorems and the tables its statement names. The one-act clause is discharged at the derived act at Part III. What stands open is the one-act clause as a premise of the persistence theorem, the two readings of the subdivision unit printed and neither chosen, and the basis label of the read a free datum at the finite identification face.
\begin{chbridge}
The literature calls the act a Bell-basis projection of two qudits, an entanglement swap (exact), the pair a maximally entangled pair, a Bell pair generalized to any two dimensions (exact), the common support the intersection of the supports of two projectors (exact), the matching a basis identification across a pair (exact), the read a projective measurement passed along the pair (nearest), and the objects named entanglement swapping \cite{Zukowski1993} and entanglement distillation concentrating noisy pairs into fewer pairs of higher fidelity \cite{Bennett1996,Nielsen2000}, by name only. The construction supplies the states, projections and the Born rule it uses at Sections~\ref{mod:FRAMEWORK} and \ref{mod:I-3-records-reads-and-the-born-count-the-gain-lemma}.
\end{chbridge}
\begin{chproof}[census and exact cell]
The premises are the Axiom and the six definitions of Section~\ref{mod:FRAMEWORK} (definitions), the axiom chapter of Section~\ref{mod:I-1-axiom-and-the-six-definitions} (an open result, read at its face), the link and its genres of Section~\ref{mod:I-2-link-and-its-six-genres}, the records and the Born count of Section~\ref{mod:I-3-records-reads-and-the-born-count-the-gain-lemma} and the moves of Section~\ref{mod:I-4-moves-in-one-table} (theorems). Nothing outside that list is used, and the chain from the axiom to this result is four chapters deep.
\emph{The exact part.} The object is the standard configuration of a genre $(d_{1}, d_{2})$. Four qudits $A$, $A'$, $B$, $B'$ at the dimensions $d_{1}, d_{1}, d_{2}, d_{2}$ with $A$ paired to $A'$ and $B$ to $B'$ at their pair kets on their common supports, a configuration at exactly its own matching among the three matchings of four qudits. From the projection of $A$ and $B$ onto the maximally entangled pair on their common support $m = \min(d_{1}, d_{2})$, by the act, we obtain the configuration $AB$ with $A'B'$ at the pair face, the squared amplitude one over the product of the dimensions,
\begin{equation}
\lvert \langle \Phi_{AB} \Phi_{A'B'} \mid \Phi_{AA'} \Phi_{BB'} \rangle \rvert^{2} = \frac{1}{d_{1} d_{2}}: \quad \tfrac{1}{4}, \ \tfrac{1}{6}, \ \tfrac{1}{8}, \ \tfrac{1}{9}, \ \tfrac{1}{12}, \ \tfrac{1}{16} \ \text{at the six genres},
\end{equation}
the squared norms of the pair kets $4, 6, 8, 9, 12, 16$ and the maximal rank of a flat pair constructible as a pure state of the two qudits equal to the common support (the ranks $1, 2$ at $(2,2)$, $(2,3)$ and $(2,4)$, $1, 2, 3$ at $(3,3)$ and $(3,4)$, $1$ to $4$ at $(4,4)$), the full product of rank the larger dimension not constructible at the three unequal genres. The smallest instance is $(2,2)$. The pair ket $(1, 0, 0, 1)$ of squared norm four, the register whole sixteen and cut nine, the projection onto it of rank one, the squared amplitude one quarter.
\emph{The census part.} The set counted is the thirteen listed acts at each of the six genres, 69 operators with their outcomes. From each act, by its output on the standard configuration, we obtain whether that output is a configuration of pairs, whether $A$ and $B$ form a pair with their partners in a pair (the pair face), the rank of the operator and the Schmidt rank of its output. The property tested, that the pair-face set equals the projections onto a maximally entangled pair on the common support, holds at 6 of 6 genres with no other candidate. A unitary on the two leaves the reduced state at the purity of the input (no pair of $A$ and $B$), the swap a configuration at another matching at the three equal genres and the entangling unitary no configuration at 6 of 6, the partial trace no configuration, a sub-maximal projection no pair, the rank equal to the Schmidt rank of the output at 69 of 69. The four steps of the chain are walked at the words of the text (onto a pair, two qudits, projecting, on their common support), derived at 4, 2, 4 and 3 cells with no posed step. The three further theorems are stated in their own boxes. The certainty clause at 120 of 120 Gram identities, the one-act clause at 32851 of 32851 instances, and the 1653 subdivisions each one read and three joins.
The check is by the counter-tests. The ontology clause struck admits every alternative at the pair face and moves 294 cells, the deciding cells to no. A non-maximal projection admitted moves 129 cells, the forced cell and the closure cells. The text struck leaves the gate unmet and moves 1173 cells. A reading offered as a premise, a hand table of acts and a network outside the fixed library are each excluded. Nine controls hold, the common supports at 1691 of 1691 links and the pair ket re-met at 20 of 20 cells among them.
The result was reproduced by an independent route, two engines written apart agreeing at every act, genre and cell.
\chcert{69}
\end{chproof}
\begin{figure}[h]
\begin{center}
\setlength{\unitlength}{1pt}
\begin{picture}(300,70)(0,0)
\put(30,50){\circle*{6}}
\put(30,10){\circle*{6}}
\put(30,50){\line(0,-1){40}}
\put(10,47){$A'$}
\put(10,7){$A$}
\put(100,50){\circle*{6}}
\put(100,10){\circle*{6}}
\put(100,50){\line(0,-1){40}}
\put(106,47){$B'$}
\put(106,7){$B$}
\put(135,28){$\longrightarrow$}
\put(200,50){\circle*{6}}
\put(270,50){\circle*{6}}
\put(200,50){\line(1,0){70}}
\put(180,47){$A'$}
\put(276,47){$B'$}
\put(200,10){\circle*{6}}
\put(270,10){\circle*{6}}
\put(200,10){\line(1,0){70}}
\put(180,7){$A$}
\put(276,7){$B$}
\end{picture}
\end{center}
\caption{The act of distilling at the standard configuration of a genre. The pairs $AA'$ and $BB'$ at the input, and at the output of the projection of $A$ and $B$ onto a maximally entangled pair on their common support the pairs $AB$ and $A'B'$, with the squared amplitude one over the product of the dimensions.}
\end{figure}

\par\medskip\noindent{\itshape Distillation read, the next chapter asks what a read costs.}\label{narr:close:13-4}\par
\section{The cost of a read and the adiabatic class}\label{mod:XIII-5-cost-of-a-read-and-the-adiabatic-class}
\chdesig{15-5}{13-5}{XIII-5-cost-of-a-read-and-the-adiabatic-class}{XIII-5-cost-of-a-read-and-the-adiabatic-class}
\label{narr:open:13-5}Theorem~234.1 reads the enumeration of the records as the cost of a read, Open result~234.2 states the class partition under one move as an open result, and Theorem~234.3 is the adiabatic theorem at the untouched records.
\begin{chtheorem}[The enumeration of the records as the cost of a read]\label{box:XIII-5-cost-of-a-read-and-the-adiabatic-class:1}
At the sixteen networks of the item and at the sixty-four library networks the records of every link are counted by the closed count law at the slot polarity by order zero to three and by kind, the orders and the kinds summing to the records at every link and network (64 of 64 and 18 of 18 defined); the record sums under the lowering at the depths zero, one and two grow with the ratio of successive sums non-increasing, the finiteness beyond the second depth open, and the growth with the order terminates at the third order at every depth; so the cost of a read at a network at an order is the record count at that order.
\end{chtheorem}
\chstanding{theorem}
\chdeps{Definitions 1.1 to 1.6, Theorems 18.1 to 18.4 and Open results 18.5 and 18.6, Theorem 230.1, Open results 241.1 and 241.2, Theorems 241.3 and 241.4 and Open results 241.5 to 241.7 and Open results 245.1 and 245.2, Theorem 245.3, Open results 245.4 to 245.8 and Theorem 245.9}{Theorems 235.1 to 235.4 and Open result 235.5}
\begin{chwall}[The class partition under one move]\label{box:XIII-5-cost-of-a-read-and-the-adiabatic-class:2}
The twenty-eight moves applied once at every library network where their form holds give 5707 applications; at the deciding reading the class partition by genre and kind is carried whole at 4124 applications and parts at 1583, the first parting the attach at the odd leg of the first library network at its qubit-ququart link, where the kept-leaf class of 834 records splits; at the local reading the invariance holds at every application, printed beside and not chosen; the adiabatic theorem by name is a gloss with its missing clause named.
\end{chwall}
\chstanding{open result}
\chdeps{Definitions 1.1 to 1.6, Theorems 18.1 to 18.4 and Open results 18.5 and 18.6, Theorem 230.1, Open results 241.1 and 241.2, Theorems 241.3 and 241.4 and Open results 241.5 to 241.7 and Open results 245.1 and 245.2, Theorem 245.3, Open results 245.4 to 245.8 and Theorem 245.9}{Theorems 235.1 to 235.4 and Open result 235.5}
\begin{chtheorem}[The adiabatic theorem at the untouched records]\label{box:XIII-5-cost-of-a-read-and-the-adiabatic-class:3}
At every one of the 5707 applications every class of the records at the links touching no move site lands whole in the class of the same name at the result, the kind of every one of the 270584 untouched links unmoved, with no parting; so the adiabatic theorem is derived at the untouched records by a sketch of five lines from the locality of the kind rule, the move action on the grade and the placement classes, the slow-change clause fenced from the sketch.
\end{chtheorem}
\chstanding{theorem}
\chdeps{Definitions 1.1 to 1.6, Theorems 18.1 to 18.4 and Open results 18.5 and 18.6, Theorem 230.1, Open results 241.1 and 241.2, Theorems 241.3 and 241.4 and Open results 241.5 to 241.7 and Open results 245.1 and 245.2, Theorem 245.3, Open results 245.4 to 245.8 and Theorem 245.9}{Theorems 235.1 to 235.4 and Open result 235.5}
The first theorem is conditional on the text. The open result has its fallback established as the theorem at the untouched records. The third theorem is conditional on the open results it rests on and on the complete list of literature names. The finiteness of the record sums beyond the second depth is named and left open.
\begin{chbridge}
The literature calls a record an invariant jointly fixed by a set of projective measurements (nearest), the slot a tensor factor of the register (exact), the link a maximally entangled pair of two qudits (exact), the move an elementary operation on the network of pairs (nearest), and the objects named the resource cost of a computation with the erasure bound of Landauer \cite{Landauer1961} and the adiabatic theorem of a slowly varying Hamiltonian \cite{BornFock1928,Berry1984}, by name only. The construction supplies the states and projections it uses at Section~\ref{mod:FRAMEWORK}.
\end{chbridge}
\begin{chproof}[census]
The premises are the Axiom and the six definitions of Section~\ref{mod:FRAMEWORK} (definitions), the records and the Born count of Section~\ref{mod:I-3-records-reads-and-the-born-count-the-gain-lemma}, the moves of Section~\ref{mod:I-4-moves-in-one-table} and the saturation of Section~\ref{mod:I-5-saturation-in-one-figure-seeds-waves-termination-sinks} (theorems), and the closure as the computation of Section~\ref{mod:XIII-1-closure-as-the-computation} (a theorem). Nothing outside that list is used, and the chain from the axiom to this result is six chapters deep.
The set counted is the records of every link of the sixteen networks of the item and of the sixty-four networks of the fixed library, and beneath them the records of the classes under the lowering at the depths one and two (23627 classes at the first depth, 8664 at the second). From a link of a genre, by the closed count law at the slot polarity, we obtain its records by order zero to three and by kind (same site, one link, cross link). At the $(2,3)$ genre the records number 172 at the orders $0\!:\!79$, $1\!:\!90$, $2\!:\!3$, and at the electron body the first link, of the genre $(3,3)$, carries
\begin{equation}
412 = 184 + 200 + 28 + 0 \quad (\text{the orders}) \ = \ 156 + 256 + 0 \quad (\text{the kinds}) ,
\end{equation}
the orders and the kinds each summing to the records. From the links of a network, by the sum, we obtain its records, and the smallest instance is the electron body at 9 sites and 10 links,
\begin{equation}
3685 = 1701 + 1854 + 130 + 0, \qquad \text{the kinds } 1669 + 2016 + 0 ,
\end{equation}
with the muon body at 5675, the tau at 7665 and the docked up quark at 344, defined at 18 of 18 item networks and 64 of 64 library networks. The order-three records are zero at every network and depth, so the growth with the order terminates at the third order, which is the first property tested. From the lowering, by the sum over the representative of every class, we obtain the record sums at the depths zero, one and two,
\begin{equation}
R_{0} = 697298, \qquad R_{1} = 26340871, \qquad R_{2} = 340474529 ,
\end{equation}
the ratios of successive sums $R_{1} / R_{0} = 26340871 / 697298$ and $R_{2} / R_{1} = 340474529 / 26340871$ non-increasing, so the growth with the depth is bounded by the first ratio at the depths computed and its finiteness beyond the second depth stands open, the two readings of the cost of the direct sum at the order two (the two-exchange records alone, 130 at the electron body, the records through the order two, 3685) are printed and neither chosen. The result is the first theorem. The cost of a read at a network at an order is the record count at that order, every count a finite exact integer. The open result and the third theorem are stated in their own boxes. The class partition under one move parting at the attach at the odd leg of the first library network, and every one of the 270584 untouched links unmoved in kind.
The check is by the counter-tests. The record clause struck, admitting records of more than two occupied slots at a link, moves the counts at 82 of 82 networks. The depth capped one step short moves the growth reading, the second ratio absent. The nucleon census re-counted at the order three equals its parent at 3857 with no differing cell. The sign reading stays the same. The text struck moves every cell. Ten controls hold, the census of the proton star re-met at 12 of 12 classes and 33 of 33 links among them.
The result was reproduced by an independent route, two engines written apart agreeing at every link, order and depth.
\chcert{23627}
\end{chproof}

\par\medskip\noindent{\itshape The Part closes on the transport of a basis and its holonomy.}\label{narr:close:13-5}\par
\section{Holonomy and the record transport}\label{mod:XIII-6-holonomy-and-the-record-transport}
\chdesig{15-6}{13-6}{XIII-6-holonomy-and-the-record-transport}{XIII-6-D542825-w-olth}
\label{narr:open:13-6}Theorem~235.1 reads the read transport from the pair, Theorem~235.2 the holonomy of the basis transport around a closed path, Theorem~235.3 the step generator as a number operator, and Theorem~235.4 the carrier phase datum at the four-cycle, four theorems, with the four open results beside the transport gathered in Open result~235.5.
\begin{chtheorem}[The read transport from the pair]\label{box:XIII-6-holonomy-and-the-record-transport:1}
The third clause of the read line is a corollary of the Axiom word maximally entangled with the first two clauses: the far set is an orthogonal basis of the far support at every genre and orientation, and the orientation sign is needed at no read.
\end{chtheorem}
\chstanding{theorem}
\chdeps{Definitions 1.1 to 1.6, Theorems 2.1 to 2.3, Open results 2.4 and 2.5, Theorem 2.6, Open results 2.7 and 2.8 and Theorems 2.9 and 2.10, Theorems 9.1 to 9.7 and Open result 9.8, Theorems 18.1 to 18.4 and Open results 18.5 and 18.6, Theorems 28.1 and 28.2, Open results 28.3 and 28.4 and Theorem 28.5 and Theorem 234.1, Open result 234.2 and Theorem 234.3}{Open result 123.1, Open result 144.1, Open result 236.1 and Open result 256.1}
\begin{chtheorem}[The holonomy of the basis transport around a closed path]\label{box:XIII-6-holonomy-and-the-record-transport:2}
At the 341 closed paths of the deciding set of the library the transport composed around the path at the identity matching is an exact unitary on the register of the starting end, defined in the ring of the twenty-fourth root of unity, nontrivial at the 333 odd paths and trivial at the 8 even ones; the Berry phase by name is the eigenphases at the identity matching, a gloss.
\end{chtheorem}
\chstanding{theorem}
\chdeps{Definitions 1.1 to 1.6, Theorems 2.1 to 2.3, Open results 2.4 and 2.5, Theorem 2.6, Open results 2.7 and 2.8 and Theorems 2.9 and 2.10, Theorems 9.1 to 9.7 and Open result 9.8, Theorems 18.1 to 18.4 and Open results 18.5 and 18.6, Theorems 28.1 and 28.2, Open results 28.3 and 28.4 and Theorem 28.5 and Theorem 234.1, Open result 234.2 and Theorem 234.3}{Open result 123.1, Open result 144.1, Open result 236.1 and Open result 256.1}
\begin{chtheorem}[The step generator as a number operator]\label{box:XIII-6-holonomy-and-the-record-transport:3}
At every one of the 1692 registers of the step-generator result the step generator equals the ring quotient multiplied by the cycle-index operator, a number operator with an equally spaced spectrum whose spacing is the quotient of that register and whose level count is the cycle order.
\end{chtheorem}
\chstanding{theorem}
\chdeps{Definitions 1.1 to 1.6, Theorems 2.1 to 2.3, Open results 2.4 and 2.5, Theorem 2.6, Open results 2.7 and 2.8 and Theorems 2.9 and 2.10, Theorems 9.1 to 9.7 and Open result 9.8, Theorems 18.1 to 18.4 and Open results 18.5 and 18.6, Theorems 28.1 and 28.2, Open results 28.3 and 28.4 and Theorem 28.5 and Theorem 234.1, Open result 234.2 and Theorem 234.3}{Open result 123.1, Open result 144.1, Open result 236.1 and Open result 256.1}
\begin{chtheorem}[The carrier phase datum at the four-cycle]\label{box:XIII-6-holonomy-and-the-record-transport:4}
At every carrier carrying a qutrit-qubit register the carrier phase datum equals the primitive fourth roots of unity in the holonomy spectrum of the transport-generator four-cycle of that register, the four-cycle present at the W, the Z and the photon words and absent at the gluon.
\end{chtheorem}
\chstanding{theorem}
\chdeps{Definitions 1.1 to 1.6, Theorems 2.1 to 2.3, Open results 2.4 and 2.5, Theorem 2.6, Open results 2.7 and 2.8 and Theorems 2.9 and 2.10, Theorems 9.1 to 9.7 and Open result 9.8, Theorems 18.1 to 18.4 and Open results 18.5 and 18.6, Theorems 28.1 and 28.2, Open results 28.3 and 28.4 and Theorem 28.5 and Theorem 234.1, Open result 234.2 and Theorem 234.3}{Open result 123.1, Open result 144.1, Open result 236.1 and Open result 256.1}
Each theorem is conditional on the open results, the conditional theorems and the tables its statement names. The Pancharatnam phase at the cap bases is placed here and read at its grade.
\begin{chwall}[The four open results beside the transport]\label{box:XIII-6-holonomy-and-the-record-transport:5}
What stands open is first the WW block at the read transport, where only the first of the four centre reads of the cap is carried to a far basis at the qutrit-qubit link, one eigenvector of six and not the five; second the record transport to the dock register, undefined as a map on the records at the electron joined network and path-dependent at the odd closed paths, the completeness relation holding neither against the identity nor at a scaling; third the carrier phase datum at the hub-pair holonomy, equal at one of four carriers with the W the first parting; and fourth the record transport at the holonomy wash, defined where the first-path transport was undefined, yet with the completeness relation parting at the washed operator too.
\end{chwall}
\chstanding{open result}
\chdeps{Definitions 1.1 to 1.6, Theorems 2.1 to 2.3, Open results 2.4 and 2.5, Theorem 2.6, Open results 2.7 and 2.8 and Theorems 2.9 and 2.10, Theorems 9.1 to 9.7 and Open result 9.8, Theorems 18.1 to 18.4 and Open results 18.5 and 18.6, Theorems 28.1 and 28.2, Open results 28.3 and 28.4 and Theorem 28.5 and Theorem 234.1, Open result 234.2 and Theorem 234.3}{Open result 123.1, Open result 144.1, Open result 236.1 and Open result 256.1}
At one face each of the four would meet its object. At the faces this chapter inherits they part as stated.
\begin{chbridge}
The literature calls the transport the passing of a basis along a pair to its far end, fixed by the pair and chosen by no connection (nearest to parallel transport along an edge), the far end the partner qudit of an entangled pair (exact), the frame a basis choice at a site (exact), the step generator the Hamiltonian that generates the step (nearest), and the objects named the holonomy of a connection with its geometric phases \cite{Berry1984,Simon1983} and a generalized measurement as a resolution of the identity \cite{Nielsen2000}, by name only. The construction supplies the wash at one act at Section~\ref{mod:II-4-measurement-the-wash-at-one-act-and-the-povm}.
\end{chbridge}
\begin{chproof}[census and exact cell]
The premises are the Axiom and the six definitions of Section~\ref{mod:FRAMEWORK} (definitions), the axiom chapter of Section~\ref{mod:I-1-axiom-and-the-six-definitions} (an open result, read at its face), the link and its genres of Section~\ref{mod:I-2-link-and-its-six-genres}, the records and the Born count of Section~\ref{mod:I-3-records-reads-and-the-born-count-the-gain-lemma} and the read at a register of Section~\ref{mod:I-6-read-at-a-register-the-probe-and-the-measurement} (theorems), and the cost of a read of Section~\ref{mod:XIII-5-cost-of-a-read-and-the-adiabatic-class} (a theorem). Nothing outside that list is used, and the chain from the axiom to this result is seven chapters deep.
\emph{The exact part.} The object is the hub triangle of the body named BA at its qubit start. The sites $f_{1}$ and $f_{2}$ of dimension two and the hub $F$ of dimension four, the closed path $f_{1} \to f_{2} \to F \to f_{1}$ along the links of the genres $(2,2)$, $(2,4)$ and $(4,2)$, odd at three links. The ring is built from its order twenty-four alone, $z$ the primitive root with the conjugate the automorphism $z \mapsto z^{-1}$, and the cap bases at the letters $Z$, $X$ and $XZ$ of the qubit are the projector columns at the roots of unity, orthogonal at one norm. From the three links, by the transport (the conjugate on the common support at every link, padded at the far register), composed around the path at the identity matching, we obtain an exact unitary on the register of dimension two of the start at every letter. At $Z$ and $X$ the identity, the spectrum $z^{0}$ twice, and at $XZ$
\begin{equation}
H_{(XZ)} = \begin{pmatrix} 1 & 0 \\ 0 & z^{12} \end{pmatrix} = \begin{pmatrix} 1 & 0 \\ 0 & -1 \end{pmatrix}, \qquad \text{spectrum } z^{0}, z^{12}, \quad \det = z^{12}, \quad \text{order } 2 ,
\end{equation}
the returned set the labels permuted ($6 \mapsto 18$, $18 \mapsto 6$ in twenty-fourths), so the path is nontrivial, the first nontrivial path of the deciding set. The spectrum is the same under every rotation of the base point and under the reverse orientation.
\emph{The census part.} The set counted is the 341 closed paths of the deciding set: 114 core triangles at the dimensions $(3,3,3)$ at the bodies with a core, 215 hub triangles at $(2,2,4)$ at the hub-bearing bodies, the five qubit cycles of four, six and eight sites of the six-root closure, and the four triangles and three squares of the four-qutrit network. From each path, by the composition above at every one of the three cap letters of its start, we obtain the holonomy, defined as an exact unitary at 341 of 341 and closed at every one. From its spectrum, by the eigenphase exponents (the exponents $0, 4, 8, 12, 16, 20$ of twenty-four turns, six distinct spectra over the set), we obtain the counts by class,
\begin{equation}
\text{nontrivial } 333 \ (\text{every odd path}), \qquad \text{trivial } 8 \ (\text{every even path}), \qquad 333 + 8 = 341 ,
\end{equation}
the odd paths carrying the half turn at the swapped directions and the sixth-root phases at the $XZ$ reads of a qutrit. The property tested, the holonomy defined at every path and nontrivial at some, holds, which is the second theorem. The reading with the matching free at every corner moves the spectrum at every one of the 341 paths and fixes it at none, so the Berry phase by name is the eigenphases at the identity matching, a gloss. The other three theorems are stated in their own boxes. The far set an orthogonal basis at every genre and orientation, the generator a number operator at 1692 registers, and the fourth-cycle datum at the carriers.
The check is by the counter-tests. The orientation of the second link of the first path inverted leaves the walk open at $f_{2}$ and moves 63 cells, the deciding cell first. The known clause struck at the first corner leaves the matching unfixed and moves 312 cells. The path opened at its last link leaves the walk not closed at $F$ and moves 66 cells. The text struck moves every cell. Thirteen controls hold, every triangle of the two kinds re-met among the cycles of the holonomy-sign table at 329 of 329 among them.
The result was reproduced by an independent route, two engines written apart agreeing at every path, letter and spectrum.
\chcert{341}
\end{chproof}
\begin{figure}[h]
\begin{center}
\setlength{\unitlength}{1pt}
\begin{picture}(220,90)(0,0)
\put(40,20){\circle*{6}}
\put(180,20){\circle*{6}}
\put(110,80){\circle{18}}
\put(40,20){\vector(1,0){134}}
\put(176,24){\vector(-1,1){58}}
\put(102,74){\vector(-1,-1){56}}
\put(28,6){$f_{1}$ (2)}
\put(170,6){$f_{2}$ (2)}
\put(122,76){$F$ (4)}
\put(98,8){$(2,2)$}
\put(150,52){$(2,4)$}
\put(50,52){$(4,2)$}
\end{picture}
\end{center}
\caption{The closed path of the hub triangle at its qubit start, $f_{1} \to f_{2} \to F \to f_{1}$. The transport composed around it at the identity matching is the unitary $\mathrm{diag}(1, -1)$ at the letter $XZ$, the first nontrivial holonomy of the deciding set.}
\end{figure}

\section{The holonomy of the transport identity at the fan loops}\label{mod:XIII-6-D542825-w-olth}
\begin{chwall}[The holonomy of the transport identity at the fan loops]\label{box:XIII-6-D542825-w-olth:1}
Face A, the flux clause closed at the stabiliser of the pair ket, reads no at its element conjunct. Face B reads yes with the first loop parting named: at the six cells of the web table the pairs of monomial matrices of fourth-root phases at the two ends fixing the pair ket number 32, 128, 1024, 384, 1536 and 6144, their projection to the support end equal to the re-matching group of the gauge-group result at 6 of 6 cells, and the transport identity of the Lorentz-group-at-the-pair result composed on the matrix units of the common support around every one of the 89607 fan loops is the identity at the 23332 even loops and the swap of the two factors at the 66275 odd loops, its determinant exponent in twenty-fourths equal to the wash integer at 89607 of 89607 loops and the lift per web equal to the gauge-group result at 45739 of 45739 webs, the twirl re-read as the average over the stabiliser at the (2,3) cell at the four read bases at the counts 0, 0, 0, 0, basis-independent; but the holonomy is an element of the group of the stabiliser at the 23332 even loops alone and outside it at the 66275 odd loops, the swap carried by no product of a change of basis with its conjugate at the dimensions 2 and 3, its determinant minus one against the determinant one of every element, the first parting at the first odd fan loop of the first subdivided web of the one (2,3) link, outside against inside; the number reading beside, the determinant exponent alone against the wash integer, stands at its own face A with the first parting none, printed and not selected.
\end{chwall}
\chstanding{open result}
\chdeps{Definitions 1.1 to 1.6, Theorem 54.1, Theorems 65.1 and 65.2 and Theorems 235.1 to 235.4 and Open result 235.5}{no theorem beyond it at this release}
The stabilizer of the pair ket is computed at every cell of the web table, the transport identity is carried around every fan loop and its holonomy printed with its determinant exponent and its membership in the group of the stabilizer, the wash integer is compared loop by loop, and the twirl is re-read as the stabilizer average. The wash integer closes at the number and the element does not, the first loop parting named. The holonomy and the transport of records stand at their own chapter.
\begin{chbridge}
In the literature the phase accumulated around a closed loop by an adiabatic transport is a holonomy of a connection on a bundle, its flux through the loop a Chern number, and the gauge transformations at a point form the structure group that the holonomy belongs to \cite{Berry1984,Wilson1974}. The correspondence is by name.
\end{chbridge}
\begin{chproof}[census]
\textit{The premises.} The axiom and the six definitions of Section~\ref{mod:FRAMEWORK}, definitions. The holonomy and the transport of records of Section~\ref{mod:XIII-6-holonomy-and-the-record-transport}, theorems and open results. The gauge group as a corollary of one construction of Section~\ref{mod:IV-2-gauge-group-as-a-corollary-of-one-construction}, a theorem. The Lorentz group and the signature of Section~\ref{mod:III-3-lorentz-group-and-the-signature}, a theorem. The chain from the axiom to this result is five chapters deep.

\textit{The census.} The set counted is the 89607 fan loops of the fan-loops result at 45739 webs of the web table, the fan counts over the webs summing as 1919 at one, 43788 at two, 16 at three and 16 at four, every walk closed and the site dimensions resolved at 89607 of 89607 with the minimal dimension re-met. At each of the six cells, (2,2) to (4,4), the stabilizer is the set of pairs of monomial matrices of fourth-root phases at the two ends fixing the pair ket, of the orders
\begin{equation}
32,\ 128,\ 1024,\ 384,\ 1536,\ 6144,
\end{equation}
the brute force over every pair of the joint group equal to the solved route at 5 of 5 cells, the determinant-one parts 32, 32, 256, 384, 384 and 6144 with their centers of orders 4, 4, 8, 4, 4 and 4, the projection to the support end equal to the re-matching group at the orders 32, 384 and 6144 at the dimensions 2, 3 and 4, the (2,3) cell re-met at the stabilizer result at 5 of 5. The property tested at every fan loop is the holonomy of the transport identity composed on the matrix units of the common support: the identity at an even link count and the swap of the two factors at an odd, the thirteen kinds of loop by minimal dimension and link count censused at 18, 21922, 3, 21940, 44155, 135, 18, 162, 1176, 15, 18, 21 and 24, the determinant exponent in twenty-fourths 0 at the identity and 12 at the swap, the census 0 at 23332 and 12 at 66275, the loop orientation reversed moving nothing at 89607 of 89607. The exponent equals the wash integer at 89607 of 89607 loops and the lift per web, the sum over the fans over 24, equals the gauge-group result at 45739 of 45739 webs, the census of the lifts 12329 at 0, 545 at one half and 32865 at 1. The holonomy is inside the group of the stabilizer at the 23332 even loops and outside at the 66275 odd loops, the swap of determinant $-1$ on the support pair space carried by no product of a change of basis with its conjugate at the dimension 2 nor 3, whose determinants read 1 at 32 of 32 and 384 of 384; the first outside loop is the first odd fan of the first subdivided web of the one (2,3) link, the holonomy $(0, 2, 1, 3)$ on the four units against the identity $(0, 1, 2, 3)$ at the first even fan, the smallest cases worked on the page. The twirl re-read as the stabilizer average at the (2,3) cell at the four read bases gives the counts 0, 0, 0, 0 from the bare counts 1, 2, 2, 2, basis-independent. So face A reads no at the element conjunct and face B reads yes with the first loop named, the number reading beside at its own face A. The result is certified at the repository.

\textit{The check.} A struck link of a loop moves the closed loops to 89606 of 89607, a (3,3) link inserted moves the first parting of the wash integer to the first fan, the transport reversed moves the equality to the even loops alone, a tree with no loop moves the fan total to 0, and the text struck excludes the run, 5 of 5 as pinned. The controls read the same at 13 of 15, the two that differ a single census string whose three kinds the engines order by value and the entries by string, the lift per web equal at 45739 of 45739 beneath it, so the control letter reads fails at the pinned grammar beside, disclosed and deciding no grade. The result was reproduced by an independent route, two engines of different construction agreeing at 0 differing keys with the runs byte-identical.
\chcert{89607}
\end{chproof}
\par\medskip\noindent{\itshape Information and computation read, Part XVI takes stock of the whole.}\label{narr:close:13-6}\par
